% Generated from validated Article Draft Result. Do not hand-edit; revise the structured draft instead.

code generationでは、研究benchmarkと利用者が日常的に触れるproductを分けて見る必要がある。AlphaCodeはprogramming problemを解く研究系のanchorであり、Codex applicationsはcode generationを実際のapplicationへ結び付ける流れを示す。その延長線上でGitHub Copilotのgeneral availabilityは、生成modelがeditor内のdeveloper workflowへ常時組み込まれるproduct boundaryとして重要である。研究上のcoding能力と、deploymentされた支援productの成熟度は同じ尺度ではない。\autocite{src-21b0842f98ec,src-8170b5e4c279,src-b1418e7ae779,src-dc0d76e03fc6}


VPTが示す方向は、生成されたtoken列を読むだけではなく、model outputをenvironment内のactionへ変換することである。ここでは成功条件が文章品質から、観測されたstateに応じて適切なaction sequenceを生成できるかへ移る。生成modelが外部環境へ作用する場合、interfaceの設計、action space、feedback loopがmodel本体とは別のsystem要素になる。\autocite{src-4e33c3e8216e,src-d117f3860a91}


ReActはreasoning traceとactionを交互に扱う構成として、推論と外部操作を一つのloopへ結び付ける方向を示した。重要なのは、Chain-of-Thought型の『考え方を出す』ことと、実際に外部actionを実行することを区別する点である。reasoning traceは内部・中間表現の扱いに関する方法であり、action executionにはtool/environmentとの接続、失敗時の観測、再計画という追加のsystem条件が生じる。\autocite{src-811e1008822e}


Atlasのようなretrieval-augmented方向では、model内部に保持されたparameterだけをknowledge sourceとせず、外部corpusから情報を取得するinterfaceが設計対象になる。これはmodel sizeを増やすこととは異なる拡張軸であり、retriever、index、取得したcontextの利用方法まで含むsystemになる。2022年には、generationを外部knowledgeへ接続することが独立したarchitecture choiceとして明確になった。\autocite{src-039e57a4916e}


code、environment action、retrievalを横断すると、共通しているのは生成結果がmodelの外側のsystemへ流れ込むことである。developer IDE、ゲーム環境、search/retrieval corpusでは、それぞれoutputの意味と失敗の影響が異なる。そのため2022年の変化は『生成能力が上がった』だけではなく、modelを外部実行・外部knowledgeへ接続するinterface layerが技術領域として立ち上がったことにある。\autocite{src-039e57a4916e,src-21b0842f98ec,src-4e33c3e8216e,src-811e1008822e,src-8170b5e4c279,src-b1418e7ae779,src-d117f3860a91,src-dc0d76e03fc6}


\begin{claimboundary}
各研究のbenchmark成績や各productの有用性は、著者・project・vendorが提示した条件に帰属する。本稿は研究problemでの結果をdeveloper productの一般的生産性へ、あるいはsimulated environmentでのactionを任意の現実世界tool実行能力へ拡張解釈しない。\autocite{src-811e1008822e,src-b1418e7ae779}
\end{claimboundary}
